\documentclass{scrreprt}

\usepackage[utf8]{inputenc}
\usepackage[francais]{babel}
\usepackage[T1]{fontenc}
\usepackage{lmodern}

\usepackage{graphicx}

\begin{document}

\thispagestyle{empty}

\phantom{~} % lolilol

\vspace{7em}

\hrule

\vspace{0.5em}

\begin{center}
	\Huge \textbf{Rapport technique} \\
	\LARGE Projet \og IUTBM\fg\ 
\end{center}

\vspace{0.5em}

\hrule

\vspace{5em}

\begin{center}
	\large
	\begin{tabular}{rl}
	 Marine & \bsc{Benoît} \\
	Florian & \bsc{Durosier} \\
	  David & \bsc{Nowinski} \\
	    Ali & \bsc{Sid} \\
	 Julien & \bsc{Voisin} \\
	 Cédric &\bsc{Zambelli}
	\end{tabular}

	\vspace{4em}

	\normalsize
	Vendredi 23 mars 2012 \\

	\vspace{2em}

	\includegraphics[width=2cm]{logoiut.eps}

	Institut Universitaire de Technologie de Belfort-Montbéliard \\
	Département Informatique
\end{center}

\newpage

\chapter*{Introduction}

Pour réaliser le projet IUTBM, nous avons choisi le langage Python, qui permet une
 programmation plus aisée des différents algorithmes. Nous avons choisi le module
 Pygame (basé sur la bibliothèque SDL) pour la réalisation de l'interface graphique.

Nous avons adopté une programmation modulaire : la classe Iutbm gère l'affichage de la
 fenêtre, et interagit avec les différents modules (menu, algorithmes) en envoyant les
 données de l'utilisateur (clics de la souris) et affichant les données renvoyées par
 l'algorithme (le dessin à l'écran).

\tableofcontents

\chapter{Spécifications}

Afin de fonctionner correctement, le projet IUTBM devra s'exécuter sur un ordinateur doté
 d'un système d'exploitation quelconque (Linux, Windows), d'un interpréteur Python 2
 (le projet a été programmé avec Python 2.7) et de la librairie Pygame (facilement
 installable avec l'utilitaire easy\_install).

Pour lancer le programme, il suffit de se placer dans le répertoire src, et d'invoquer
 l'interpréteur Python avec en argument le fichier iutbm.py.

Par exemple, sous le système d'exploitation Linux :

\begin{verbatim}
$ cd src/
src $ python iutbm.py
\end{verbatim}

L'interface du jeu s'ouvre alors.

\chapter{Description des fichiers}

\section{iutbm.py : la classe iutbm}

Dans le constructeur de la classe iutbm, la fenêtre principale est créée et le menu est
 initialisé avec les différentes options. La méthode main représente la boucle principale
 du programme : les différents évènements (clics de la souris, appui sur les touches clés
 du programme, fermeture de la fenêtre) sont détectés et déclenchent soit le passage
 dans un autre mode (mode algorithme, mode menu, mode aide), soit la transmission
 de l'évènement à l'algorithme en cours d'exécution.

Après la gestion des évènements, on appelle les différentes méthodes de dessin :

\begin{description}
	\item[dans le mode algorithme], on appelle la méthode \_drawInAlgo(), qui dessine
	 d'abord les trois boutons (Retour, Aide et Solution) avant d'appeler la méthode
	 \_draw() de l'algorithme en cours d'exécution, qui va lui-même dessiner le reste;
	\item[dans le mode menu], on appelle directement la méthode draw() du menu, qui
	 va dessiner le menu dans la fenêtre;
	\item[dans le mode aide], on dessine le bouton Retour puis on appelle la méthode
	 \_help() de l'algorithme en cours d'exécution, qui va dessiner le reste de la fenêtre
	 d'aide.
\end{description}

\section{algo.py : la classe abstraite des algorithmes}

La classe Algo représente un algorithme. Elle est étendue par tous les algorithmes
 présents dans le programme.

Le constructeur initialise quelques données par défaut : l'affichage initial pour le
 dessin avec Pygame, le titre de l'algorithme et sa description pour l'écran d'aide,
 et un repère standard (par défaut sur 0-100) qui servira à adapter l'affichage à la
 résolution d'écran choisie par l'utilisateur.

La méthode \_help() permet d'afficher un écran d'aide à l'utilisateur. Il affiche
 le titre de l'algorithme, les différentes lignes de la description (séparées par
 des \#), ainsi que des images représentant les boutons utilisables par l'utilisateur.

La méthode \_get\_corres\_pixel() permet de convertir les coordonnées x, y passées
 en paramètre et faisant partie du repère défini dans le constructeur, en coordonnées
 dans la fenêtre. Cela permet, sans toucher au code source de chaque algorithme, de
 pouvoir changer la résolution du jeu et d'adapter les dessins en conséquence.

La méthode \_reset() permet de réinitialiser l'algorithme.

La méthode \_draw() utilise les principales méthodes de Pygame pour afficher une
 représentation à l'écran de l'algorithme en cours.

La méthode \_update() est appelée lorsque l'utilisateur clique dans la fenêtre de jeu.
 Les coordonnées (x, y) sont passées en paramètre, ainsi que le bouton qui a été utilisé.
 Cette méthode doit permettre à l'utilisateur d'interagir avec l'algorithme, et utilise
 notamment la méthode Rect.collidepoint() pour vérifier si l'utilisateur a cliqué sur
 un élément du décor.

\section{buttonMenu.py : le menu principal}

Le menu principal est basé sur des boutons orange cliquables. Pour des raisons pratiques,
 il réutilise le système de repère implémenté dans les algorithmes. La grille est découpée
 en 6 sections dans lequelles sont placés les boutons dans la méthode draw().

La méthode update() vérifie si l'utilisateur a cliqué sur un bouton. Si c'est le cas, elle
 quitte l'application ou renvoie l'algorithme correspondant. Sinon, elle renvoie None.

\section{rotatingMenu.py : rotation du menu}

Ce fichier permet d'avoir un menu rotatif, qui présente les différents jeux a l'utilisateur.
De plus, il permet, notemment grâce a la méthodes addItem(item) et à la classe MenuItem, de générer
 une option supplémentaire du menu.

Cependant, c'est principalement une classe qui gère l'esthétique du menu, et le placement des items.

\section{Classes des jeux}
	\subsection{couplage.py : le jeu du couplage}
	Le jeu du couplage permet donc à l'utilisateur de résoudre un problème de graphe, qui en apparence est
	 un jeu intéractif. Le problème étant de pouvoir atteindre un flot maximal dans un graphe composé de différents
	 sommets, dont un sommet Source qui est le fournisseur du flot, et un sommet Puit qui est le recepteur. Les autres sommets
	 qui composent sont soitreprésentatif de personnes (les sommets liés au sommet Source), soit de pizza (l'objet convoité, les 
	 sommets liés au sommet Puit).

	Dans le fichier \emph{couplage.py}, nous avons la classe 

\end{document}
